%%%%%% Matching
\question\MatchQuestion{Heterochrony}{Changes in the timing and use of developmental genes.}
\question\MatchQuestion{Allometry}{Different parts of the body develop at different rates at different times.}
\question\MatchQuestion{Paedomorphy}{Retention of child-like traits in adult organism.}
\question\MatchQuestion{Mutualism}{Both species get ecological benefit over the long term.}
\question\MatchQuestion{Parasitism}{One species get ecologicals benefit over the long term, the other species is harmed.}
\question\MatchQuestion{Commensalism}{One species get ecological benefit over the long term, the other neither benefits or loses.}
\question\MatchQuestion{Diversity}{The number and relative abundances of species in the community.} 
\question\MatchQuestion{Richness}{The total number of species present in the community.} 
\question\MatchQuestion{Cryptic Species}{May be overlooked by the morphological species concept.}
\question\MatchQuestion{Fundamental Niche}{Based on the resources potentially used by a species.}
\question\MatchQuestion{Realized Niche}{Based on the resources actually used by a species.}
\question\MatchQuestion{Taxonomy}{The science of assigning scientific names to species.}
\question\MatchQuestion{Binomial Nomenclature}{The species name has two parts.}
\question\MatchQuestion{Prezygotic Barrier}{Different reproductive behaviors prevent copulation.}
\question\MatchQuestion{Postzygotic Barrier}{Different reproductive behaviors prevent the zygote from developing after fertilization.}
\question\MatchQuestion{Ecosystem}{Interaction between a community and the physical environment.}
\question\MatchQuestion{Primary Consumers}{Eaten by secondary consumers.}
\question\MatchQuestion{Primary Consumers}{Eat primary producers.}
\question\MatchQuestion{Secondary Consumers}{Eat primary consumers.}
\question\MatchQuestion{Secondary Consumers}{Eaten by tertiary consumers.}
\question\MatchQuestion{Primary Producers}{Eaten by primary consumers.}
\question\MatchQuestion{Primary Producers}{Responsible for gross primary production.}
\question\MatchQuestion{Mimicry}{A tasty species looks like a poisonous species.}
\question\MatchQuestion{Aposomatic Coloration}{A poisonous species warns a potential predator.}

\question\MatchQuestion{Camouflage}{A species blends in with its environment.}
\question\MatchQuestion{Competition}{Occurs when individuals require the same limited resource.}
\question\MatchQuestion{Intraspecific Competition}{Occurs when individuals of the same species require the same limited resource.}
\question\MatchQuestion{Interspecific Competition}{Occurs when individuals of different species require the same limited resource.}
\question\MatchQuestion{Allopatric}{New species evolve in geographic isolation.}
\question\MatchQuestion{Sympatric}{New species evolve without geographic isolation.}
\question\MatchQuestion{Gross Primary Production}{Total amount of sun's energy converted to organic matter.}
\question\MatchQuestion{Net Primary Production}{Amount of sun's energy transferred to primary consumers.}
\question\MatchQuestion{Trophic Cascade}{Upper trophic level species causes alternating abundance levels of lower trophic level species.}
\question\MatchQuestion{Resource Partitioning}{Species divide shared resource, increasing diversity.}
\question\MatchQuestion{Keystone Species}{A predator species that increases community diversity.}
\question\MatchQuestion{Logistic Growth}{Occurs when a population has a carrying capacity.}
\question\MatchQuestion{Exponential Growth}{Occurs when a population has unlimited resources.}


%%%%%% True False
\question\TFQuestion{T}{The best classification system is that which most closely reflects evolutionary history.}
\question\TFQuestion{T}{Mutualism can be thought of as two species that ecologically exploit each other over the long term.}
\question\TFQuestion{T}{Richness measures only the number of species in a community. }
\question\TFQuestion{T}{Evenness measures the distribution of relative abundance across all species in a community.}
\question\TFQuestion{F}{Evenness measures the distribution of relative abundance across all species in a population.}
\question\TFQuestion{T}{Net primary production represents the amount of energy that will be available to consumers in an ecosystem.}
\question\TFQuestion{F}{Net primary production represents the amount of energy that will be available to secondary consumers in an ecosystem.}

\question\TFQuestion{T}{The brightly colored pattern of poison dart frogs is an example of aposomatic coloration. }
\question\TFQuestion{T}{The brightly colored pattern of bees and wasps is an example of aposomatic coloration. }
\question\TFQuestion{F}{The brightly colored pattern of harmless coral reef fishes is an example of aposomatic coloration. }
\question\TFQuestion{T}{One factor that determines if a population is getting smaller is an increasing mortality rate.}
\question\TFQuestion{T}{One factor that determines if a population is getting larger is an decreasing mortality rate.}
\question\TFQuestion{T}{One factor that determines if a population is getting larger is an increasing birth rate.}
\question\TFQuestion{T}{One factor that determines if a population is getting smaller is an decreasing rate.}
\question\TFQuestion{T}{One factor that determines if a population is getting smaller is an increasing emigration rate.}
\question\TFQuestion{F}{A population will always grow smaller if emigration is greater than birth rate.}
\question\TFQuestion{F}{A population will always grow smaller if immigration is greater than birth rate.}
\question\TFQuestion{F}{A population will always grow smaller if immigration is greater than death rate.}
\question\TFQuestion{T}{Competitive exclusion occurs when one species outcompetes another species that uses the realized same niche.}
\question\TFQuestion{F}{If negative feedback occurs when \textit{N} is greater than \textit{K}, then positive feedback occurs when \textit{N} is less than \textit{K}.}
\question\TFQuestion{F}{Negative feedback most often occurs in a population when \textit{N} is much less than \textit{K}.}
\question\TFQuestion{T}{Negative feedback most often occurs in a population when \textit{N} is much larger than \textit{K}.}

\question\TFQuestion{T}{Disease killed off most of the sea urchins around Jamaican coral reefs.  The number of species on the coral reef decreased dramatically after the sea urchins died. The sea urchins may have been a keystone species.}
\question\TFQuestion{F}{Any one trophic level supports on average about 10\% of the biomass of the next higher trophic level.}
\question\TFQuestion{F}{Any one trophic level supports on average about 10\% of the energy of the next higher trophic level.}

\question\TFQuestion{T}{Any one trophic level supports on average about 10\% of the biomass of the next lower trophic level.}
\question\TFQuestion{T}{Any one trophic level supports on average about 10\% of the energy of the next lower trophic level.}


\question\TFQuestion{F}{A field containing 30,000 kg of plant biomass will support about 300 kg of tertiary consumer biomass.}
\question\TFQuestion{F}{The carrying capacity of source habitats decreases when the carrying capacity of sink patches increases.}
\question\TFQuestion{T}{The carrying capacity of a habitat may change as environmental conditions change.}
\question\TFQuestion{F}{Evenness measures the distribution of relative abundance across all species in a population.}
\question\TFQuestion{T}{The brightly colored pattern of bees and wasps is an example of aposomatic coloration. }
\question\TFQuestion{F}{Heterochrony can be defined as the retention of juvenile characteristics in the adult.}
\question\TFQuestion{T}{Mutualism can be thought of as two species that ecologically exploit each other over the long term.}
\question\TFQuestion{F}{Competitive exclusion will occur when two species use even one identical resource out of many resources.}
\question\TFQuestion{F}{If populations of a single species are geographically isolated for \textit{any} period of time, then they will evolve biological reproductive isolation.}
\question\TFQuestion{F}{A clade is a term used for any non-monophyletic group.}
\question\TFQuestion{F}{An insect that has evolved to resemble a plant twig will probably be able to avoid parasitism.}
\question\TFQuestion{T}{An insect that has evolved to resemble a plant twig will have a better chance of avoiding predation.}
\question\TFQuestion{T}{Energy that enters into an ecosystem is eventually lost as metabolic heat.}